Neste capítulo, exploraremos as principais características de enchentes e as suas causas. Exploraremos em especial o
tipo de enchentes que afetou o estado do Rio Grande do Sul no começo de 2024, as Enchentes Fluviais. Será explorada a
morfologia de rios que leva estas situações, bem como a intrinseca Vulnerabilidade urbana à estes eventos climáticos
extremos.

% IPCC AR6 WGI - Chapter 8
\subsection{Enchentes e suas características}
Enchentes são caracterizadas pelo aumento do nível da água em áreas que normalmente permanecem secas. Esse acúmulo
prolongado ocorre devido à dificuldade de escoamento, seja pela incapacidade do solo de absorvê-la ou pela lentidão na
drenagem natural, como o fluxo em grandes corpos d'água. Em geral, elas podem ser classificadas em três tipos
principais: enchentes pluviais, resultantes de chuvas intensas que excedem a capacidade de absorção do solo, causando
alagamentos; enchentes fluviais, provocadas pelo transbordamento de corpos d'água – geralmente rios – quando o nível da
água ultrapassa suas margens, frequentemente devido a chuvas persistentes, derretimento de gelo ou acúmulo de
sedimentos; e enchentes costeiras, que ocorrem em áreas litorâneas devido ao aumento do nível do mar, grandes ondas ou
fenômenos como ciclones e tsunamis.

As causas de uma enchente nem sempre podem ser facilmente atribuídas a um único evento específico, pois geralmente
dependem de uma combinação de fatores. No caso das enchentes fluviais, por exemplo, não é apenas a ocorrência de chuvas
intensas no corpo d'água principal ou em seus tributários que provoca o transbordamento. Se o corpo hídrico tiver
capacidade de escoar a água de forma eficiente, o alagamento pode ser evitado. Muitos corpos d'água, cuja principal
fonte de abastecimento são enchentes pluviais, possuem uma morfologia adaptada para lidar com essa condição. Outros
fatores que contribuem para o transbordamento incluem a área de escoamento disponível em canais naturais e artificiais,
a umidade prévia do solo, sua capacidade de percolação, a presença de ventos fortes e sua interação com o fluxo da água,
a taxa de evaporação após o evento, a capacidade de absorção da vegetação local, o declive do terreno em relação à área
de escoamento, entre outros. Assim, não se pode associar exclusivamente o aumento das chuvas intensas (\autoref{rain})
ao crescimento de todos os eventos extremos desse tipo (REFERÊNCIA).

% ######################################################################################################################
%EFEITOS GEOGRAFICOS QUE LEVAM A ENCHENTES -> MOFOLOGIA DE RIOS, ESCOAMENTO, INTERAÇÃO COM VENTOS, DECLIVE E CANAIS:

% ######################################################################################################################
Mesmo assim, as enchentes são um fenômeno natural e frequentemente característico de certas regiões. Áreas com invernos
rigorosos, marcados pela acumulação de neve, costumam enfrentar enchentes durante o verão, quando o rápido derretimento
do gelo transporta grandes volumes de água de uma região para outra, sustentando ecossistemas que dependem desse influxo
sazonal. Regiões influenciadas por monções — ventos sazonais que trazem para o interior um acúmulo de umidade e
precipitação marítima — também são marcadas por cheias em seus corpos d'água, essenciais para a agricultura local, como
ocorre nos tributários dos rios Ganges e Amarelo. De maneira semelhante, as cheias anuais regulares do Rio Nilo são
fundamentais para a fertilidade da região. 

Entretanto, mesmo sendo uma parte integral do ciclo da água, a presença de assentamentos humanos geralmente corresponde
com a existência de grandes corpos d'água que podem ser, dadas circunstâncias, vulneráveis à enchentes. Essas imundações
podem invadir esses locais de forma inesperada e veloz, colocando vidas em risco, comprometendo a segurança das
comunidades e causando danos estruturais que amplificam a disrupção das atividades humanas. Fatalidades causadas por
afogamentos são comuns em regiões alagadas, não apenas pela rapidez com que as enchentes ocorrem, mas também devido às
características do escoamento das águas, que podem gerar correntes fortes e imprevisíveis. As enchentes também criam
condições propícias para o agravamento de outras condições sanitárias. Elas podem transmitir doenças por meio da água
contaminada, além de comprometer sistemas críticos, como a contaminação de fontes de água potável, a interrupção do
fornecimento de energia elétrica, afetando hospitais e serviços essenciais, e a destruição de cadeias de suprimentos.
Esses fatores podem isolar comunidades, mesmo aquelas não diretamente afetadas pela inundação, dificultando o acesso a
mantimentos e recursos básicos.

Os efeitos das enchentes, no entanto, não se limitam às comunidades urbanas. Atividades econômicas fundamentais, como a
agricultura, estão entre as mais prejudicadas por esses desastres. A necessidade de grandes extensões de terra e a
proximidade com corpos d’água para irrigação tornam o setor agrícola particularmente vulnerável. Em situações de
enchente, a safra pode ser destruída não somente pela força da água, mas também pela permanência prolongada de água
nessas regiões, já que muitas espécies de plantas não conseguem sobreviver a tais condições. Além da destruição de
colheitas e da perda de gado, as enchentes frequentemente resultam em danos à infraestrutura essencial, como sistemas de
irrigação, equipamentos, instalações de processamento e transporte, além de alterações nas terras cultiváveis. De acordo
com a Organização de Alimentos e Agricultura das Nações Unidas \cite{onu}, enchentes são responsáveis por mais da metade
de todas as perdas em colheitas causadas por desastres naturais. Entre 2008 e 2018, essas perdas em países em
desenvolvimento ultrapassaram 21 bilhões de dólares, tornando as enchentes a segunda maior causa de danos econômicos no
setor agrícola, atrás apenas das secas. Em países em desenvolvimento, onde setores primários têm um peso significativo
na economia, os efeitos desses desastres são amplamente sentidos de forma indireta no enfraquecimento econômico, podendo
levar a uma maior situação de vulnerabilidade e despreparo.

\subsection{Tendências de Enchentes devido à mudanças climáticas}
De maneira direta e indireta, as atividades humanas alteram profundamente o ciclo hidrológico. Conforme demonstrado na
\autoref{rain}, o aquecimento global não apenas eleva as temperaturas, mas também intensifica a frequência e a magnitude
das chuvas, impulsionado pela maior quantidade de umidade presente na atmosfera. Embora as variações naturais possam
dificultar a identificação de mudanças locais nos eventos de chuva, há evidências de que, quando padrões meteorológicos
naturais promovem chuvas prolongadas em um clima aquecido, a intensidade desses episódios é amplificada pela abundância
de vapor d’água. Isso resulta em um aumento de eventos extremos, mesmo que, regionalmente, as mudanças possam se
manifestar de formas diversas – com áreas áridas tornando-se ainda mais secas e regiões úmidas experimentando chuvas
mais intensas.

Entretanto, é importante destacar que chuvas mais intensas nem sempre se traduzem automaticamente em enchentes mais
graves, pois os impactos dependem de múltiplos fatores, como as características da bacia hidrográfica, a topografia, a
duração do evento e o estado de umidade do solo antes da ocorrência da chuva. Em algumas regiões, por exemplo, chuvas
menos frequentes, mas extremamente intensas, podem resultar em solos secos e endurecidos, os quais apresentam menor
capacidade de absorção, favorecendo o escoamento superficial e, consequentemente, inundações em rios, lagos e áreas
rebaixadas. Além disso, a gestão do uso do solo – como o desmatamento para expansão agrícola ou a urbanização
desordenada – pode acelerar o fluxo de água, direcionando-a de forma mais rápida para áreas vulneráveis. Em
contrapartida, a extração intensiva de água dos rios pode, em certos contextos, reduzir os níveis dos cursos d’água e,
temporariamente, mitigar o risco de enchentes.

Portanto, apesar das variações naturais que tornam desafiador atribuir mudanças locais exclusivamente ao aquecimento
global, os registros indicam que o aumento da concentração de gases de efeito estufa, resultante de atividades humanas,
tem contribuído para chuvas diárias cada vez mais intensas e eventos extremos. Esse cenário, aliado às alterações no
manejo do solo e das bacias hidrográficas, aponta para uma tendência de enchentes mais severas em muitas regiões,
exigindo estratégias de adaptação e planejamento urbano para mitigar os impactos futuros.

% These natural variations make it difficult to determine whether heavy rainfall events are changing locally as a result
% of global warming. However, when natural weather patterns bring heavy and prolonged rainfall in a warmer climate, the
% intensity is increased by the larger amount of moisture in the air. An increased intensity and frequency of
% record-breaking daily rainfall has been detected for much of the land surface where good observational records exist,
% and this can only be explained by human-caused increases in atmospheric greenhouse gas concentrations. Heavy rainfall
% is also projected to become more intense in the future for most places. So, where unusually wet weather events or
% seasons occur, the rainfall amounts are expected to be greater in the future, contributing to more severe flooding.

% However, heavier rainfall does not always lead to greater flooding. This is because flooding also depends upon the
% type of river basin, the surface landscape, the extent and duration of the rainfall, and how wet the ground is before
% the rainfall event (FAQ 8.2, Figure 1). Some regions will experience a drying in the soil as the climate warms,
% particularly in subtropical climates, which could make floods from a rainfall event less probable because the ground
% can potentially soak up more of the rain. On the other hand, less frequent but more intense downpours can lead to dry,
% hard ground that is less able to soak up heavy rainfall when it does occur, resulting in more runoff into lakes,
% rivers and hollows. 

% Flooding is also affected by changes in the management of the land and river systems. For example, clearing forests
% for agriculture or building cities can make rainwater flow more rapidly into rivers or low-lying areas. On the other
% hand, increased extraction of water from rivers can reduce water levels and the likelihood of flooding.

% River flood: Emerging literature in the region documents ongoing changes in river floods. Mernild et al. (2018). In
% Brazil, floods are becoming more frequent and intense in wet regions but less frequent and intense in drier regions
% (Bartiko et al., 2019; Borges de Amorim and Chaffe, 2019), with higher propagation of hydrological changes through
% anthropogenically modified agricultural basins (Chagas and Chaffe, 2018). Record, catastrophic, unprecedented, and
% once-in-a-century flooding events have also been reported in recent decades in the tributaries of the Amazon River or
% along its mainstream (Sena et al., 2012; Espinoza et al., 2013; Marengo et al., 2013; Filizola et al., 2014), in
% Argentinean rural and urban areas (Barros et al., 2015). In the Amazon basin, the significant increase in extreme flow
% is associated with the strengthening of the Walker circulation (Barichivich et al., 2018). Available projections for
% the region show increases in river floods in SES and SAM (medium confidence). Projections indicate that SES and the
% coasts of Ecuador and Peru will experience a tendency towards wetter conditions that can be a proxy for longer periods
% of flooding and enhanced river discharges (Zaninelli et al., 2019). CORDEX models project the strongest changes for
% the peak flow with a return period of 100 years in SES by mid-century and under RCP8.5 (Figure 12.8).

\subsection{Características de Enchentes}
\lipsum[1-10]

\subsection{Impactos Ambientais}
\lipsum[1-10]

\subsection{Impactos Sociais}
\lipsum[1-10]

\subsection{Vulnerabilidade urbana à enchentes}
\lipsum[1-5]

% Começar falando que cidades geralmente sofrem mais com os efeitos das mudanças climáticas por causa de fatores como as
% ilhas de calor urbanas;

% Usos artificiais para enchentes -> Geração de energia, represas e defesa (Dutch Waterline, Stelling van Amsterdam);

% Falar aqui de trabalhos como o do Ta;

% El ÑINO:
% https://wmo.int/media/news/global-temperature-record-streak-continues-climate-change-makes-heatwaves-more-extreme